(a) Take $X=\mathbb P^1_k$ and $\mathcal E=\mathcal O_X(-1)^{\oplus2}$. By (7.11), $\Gamma(\mathbb P(\mathcal E),\mathcal O(1))=\Gamma(X,\mathcal E)=0$. A sheaf very ample relative to $X$ is generated by the pullbacks of the homogeneous coordinates of some $\mathbb P_X^N$. Hence $\mathcal O_{\mathbb P(\mathcal E)}(1)$ is not very ample relative to $X$.

(b) First, some power $\mathcal L^r$ is very ample relative to $Y$. Indeed, the proof of (7.6) applies with affine opens of $X$ contained in inverse images of affine opens of $Y$. Choose finitely many affine opens $X_{s_i}$ of this form covering $X$, with the $s_i$ sections of a common power of $\mathcal L$. Finite type gives finitely many algebra generators for each $\Gamma(X_{s_i},\mathcal O_X)$ over the coordinate ring of the corresponding affine open of $Y$. By (5.14), multiplying these generators by sufficiently large powers of $s_i$ extends them to global sections. These sections, together with the powers of the $s_i$, give an immersion into $\mathbb P_Y^N$, exactly as in (7.6). Tensoring $\mathcal L^r$ with the globally generated sheaf $\mathcal L^{q-r}$ and using the graph and Segre construction of \exref{II.7.5}(d) shows that $\mathcal L^q$ is very ample relative to $Y$ for every sufficiently large $q$.

Choose $a$ such that $\mathcal S_1\otimes\mathcal L^a$ is generated by finitely many global sections. By (7.10), $\mathcal O_P(1)\otimes\pi^*\mathcal L^a$ is very ample relative to $X$. For all sufficiently large $n$, $\mathcal L^{n-a}$ is very ample relative to $Y$, so \exref{II.5.12}(b) makes
\[(\mathcal O_P(1)\otimes\pi^*\mathcal L^a)\otimes\pi^*\mathcal L^{n-a}=\mathcal O_P(1)\otimes\pi^*\mathcal L^n\]
very ample relative to $Y$.
